\Page{Cours\quad /\quad 1}{Sources et spectres}{c-sources}
La lumière émise par un objet peut comporter une seule fréquence dans un modèle idéal, quelques raies ou une gamme continue de fréquences. Un \textbf{spectre} décrit la répartition de l'intensité lumineuse en fonction de la fréquence ou de la longueur d'onde dans le vide.
\begin{loi}{Deux idéalisations utiles}
Une source \textbf{ponctuelle} a des dimensions négligeables à l'échelle de l'expérience. Une source \textbf{monochromatique} est décrite par une seule fréquence $\nu$ et une longueur d'onde dans le vide $\lambda_0$ :
\[\result{\lambda_0=\frac{c}{\nu}},\qquad c=299\,792\,458\ \mathrm{m\,s^{-1}}.\]
Une source étendue se modélise par un ensemble de sources ponctuelles. Une source polychromatique se décompose en composantes monochromatiques.
\end{loi}
\fig{spectres}{Deux types de spectres ; les hauteurs de raies ne sont pas des données expérimentales.}
\Sub{Couleur et longueur d'onde}
Le visible s'étend approximativement de $400$ à $800\ \mathrm{nm}$, selon la sensibilité du détecteur humain. Repères : violet vers $400$ nm, bleu vers $450$ nm, vert vers $550$ nm, rouge vers $650$ nm. Ces limites ne sont pas des frontières absolues. Une couleur perçue peut aussi résulter d'un mélange spectral.
\begin{methode}{Exemple à refaire}
Pour $\nu=5{,}45\times10^{14}\ \mathrm{Hz}$, on trouve $\lambda_0\simeq550\ \mathrm{nm}$ : cette composante est verte. Le nanomètre vaut $10^{-9}\ \mathrm{m}$.
\end{methode}
\source{Modèles usuels : filament chaud, spectre continu ; lampe spectrale, raies ; laser, bande étroite que l'on assimile souvent à une composante monochromatique.}

\Page{Cours\quad /\quad 2}{Le modèle des rayons}{c-modele}
Un \textbf{rayon lumineux} représente localement une direction et un sens de propagation de la lumière. Il ne représente ni une particule matérialisée, ni un faisceau de diamètre strictement nul que l'on pourrait isoler au laboratoire.
\begin{loi}{Milieu transparent et indice}
Dans le domaine considéré, un milieu transparent laisse se propager la lumière avec une absorption négligeable. Il est homogène si ses propriétés sont indépendantes de la position et isotrope si elles sont indépendantes de la direction.
\[\result{n(\lambda_0)=\frac{c}{v}},\qquad
\lambda=\frac{v}{\nu}=\frac{\lambda_0}{n}.\]
L'indice est sans dimension ; $v$ est la vitesse de phase dans le milieu. Pour les milieux transparents ordinaires étudiés ici, $n\geq1$. La fréquence est conservée au passage d'une interface immobile.
\end{loi}
\fig{modele}{Propagation rectiligne et limite liée à la diffraction.}
\Sub{Domaine de validité}
Dans un milieu homogène et isotrope, les rayons sont rectilignes. Dans un système réciproque, le chemin peut être parcouru en sens inverse : c'est le \textbf{retour inverse de la lumière}. Les lois géométriques s'appliquent séparément à chaque longueur d'onde si le milieu est dispersif.
\begin{piege}{Une échelle doit être comparée à une autre}
L'approximation géométrique demande que les dimensions pertinentes, par exemple une ouverture, soient grandes devant $\lambda$. Elle ne décrit pas correctement la diffraction, les interférences, la polarisation ni les détails d'une propagation guidée à l'échelle de la longueur d'onde.
\end{piege}

\Page{Cours\quad /\quad 3}{Réflexion et réfraction}{c-lois}
Un \textbf{dioptre} est une interface entre deux milieux transparents. Le point d'arrivée est le point d'incidence $I$. La normale est perpendiculaire à la surface locale ; elle définit avec le rayon incident le \textbf{plan d'incidence}.
\fig{descartes}{Les angles d'incidence et de réfraction se mesurent depuis la normale.}
\begin{loi}{Lois de Snell-Descartes}
Les rayons incident, réfléchi et réfracté appartiennent au plan d'incidence. Avec des angles géométriques positifs entre $0$ et $\pi/2$ :
\[\result{i_{\mathrm{réfléchi}}=i},\qquad
\result{n_1\sin i=n_2\sin r}.\]
Les rayons incident et réfléchi se trouvent de part et d'autre de la normale. Lorsque $n_2>n_1$, le rayon réfracté se rapproche de la normale : $r<i$.
\end{loi}
\begin{methode}{Traiter un dioptre sans erreur d'angle}
Tracer la surface et sa normale ; identifier le milieu de départ ; écrire la loi avant tout calcul ; vérifier que l'argument de $\arcsin$ appartient à $[0,1]$ ; comparer $i$ et $r$ avec les indices.
\end{methode}
\Sub{Application numérique}
De l'air ($n_1=1{,}00$) vers un verre ($n_2=1{,}50$), à $i=40{,}0^\circ$ :
\[r=\arcsin\!\left(\frac{1{,}00}{1{,}50}\sin40{,}0^\circ\right)=25{,}4^\circ.\]
Le modèle fournit une direction. Il ne donne pas la proportion d'énergie réfléchie et transmise ; les coefficients de Fresnel ne sont pas nécessaires ici.

\Page{Cours\quad /\quad 4}{La réflexion totale}{c-totale}
\Sub{Une condition d'existence du rayon transmis}
La relation $\sin r=(n_1/n_2)\sin i$ n'admet un angle réel que si $(n_1/n_2)\sin i\leq1$. Si $n_1\leq n_2$, cette condition est toujours satisfaite pour un angle d'incidence admissible.
\begin{loi}{Seuil de réflexion totale}
Pour un passage d'un milieu d'indice $n_1$ vers un milieu d'indice $n_2<n_1$, on définit
\[\result{i_c=\arcsin\!\left(\frac{n_2}{n_1}\right)}.\]
Si $i<i_c$, un rayon transmis existe. Si $i=i_c$, la direction transmise est rasante. Si $i>i_c$, il n'existe plus de rayon propagatif réfracté : il y a \textbf{réflexion totale} dans ce modèle.
\end{loi}
\fig{totale}{La réflexion totale exige à la fois $n_1>n_2$ et une incidence suffisante.}
\begin{methode}{Démonstration à savoir refaire}
Partir de $\sin r\leq1$, isoler $\sin i\leq n_2/n_1$ puis utiliser la croissance de $\sin$ sur $[0,\pi/2]$. Il faut annoncer le sens du passage avant d'introduire $i_c$.
\end{methode}
\Sub{Exemple et interprétation}
Pour verre-air, $n_1=1{,}50$ et $n_2=1{,}00$, donc $i_c=41{,}8^\circ$. Un rayon dans le verre arrivant à $50^\circ$ est totalement réfléchi. Le même angle pour un rayon arrivant depuis l'air ne produit aucune réflexion totale.
\begin{piege}{Le mot « totale » concerne ici l'absence de transmission propagative}
On néglige l'absorption et on suppose une interface idéale entre milieux transparents. La description ondulatoire d'un champ évanescent dépasse ce chapitre.
\end{piege}

\Page{Cours\quad /\quad 5}{Objet et image dans un miroir}{c-miroir}
Un point objet est \textbf{réel} pour un système lorsque les rayons incidents en proviennent effectivement. Un point image est \textbf{réel} si les rayons émergents y convergent ; il est \textbf{virtuel} si seuls leurs prolongements en arrière se coupent. Un objet virtuel correspond à un faisceau incident qui convergerait vers ce point en l'absence du système.
\fig{miroir}{L'image virtuelle $A'$ est le symétrique de $A$ par rapport au miroir.}
\begin{loi}{Miroir plan et stigmatisme}
Un miroir plan associe à un point objet un point image symétrique par rapport à son plan. Tous les rayons réfléchis ont des prolongements passant par ce point : le miroir plan réalise un \textbf{stigmatisme rigoureux} dans le modèle géométrique.
\end{loi}
\Sub{Justifier la construction}
Projeter $A$ orthogonalement sur le miroir en $H$. Construire $A'$ de l'autre côté avec $HA'=HA$. Pour un point d'incidence quelconque, les triangles rectangles $AHI$ et $A'HI$ ont mêmes dimensions ; leurs angles vérifient donc la loi de réflexion. Le rayon réfléchi est aligné avec $A'I$.
\begin{piege}{Une image virtuelle est observable}
Les rayons réfléchis entrent réellement dans l'œil ou l'appareil photographique. Un instrument peut ensuite en former une image réelle sur son propre capteur. En revanche, un écran placé au point virtuel ne reçoit pas des rayons qui y convergent.
\end{piege}
Pour un objet étendu, on construit l'image de plusieurs points. Les longueurs sont conservées. La formulation usuelle « le miroir inverse droite et gauche » est imprécise : la symétrie agit suivant la direction normale au miroir.

\Page{Cours\quad /\quad 6}{Netteté et conditions de Gauss}{c-gauss}
Un système est \textbf{stigmatique} pour un couple de points si tous les rayons issus du point objet et traversant le système passent par un même point image, ou semblent en provenir. Une lentille réelle ne possède généralement pas cette propriété pour tous les rayons.
\begin{loi}{L'approximation de Gauss}
Pour un système centré, on ne conserve que les rayons \textbf{proches de l'axe optique} et \textbf{faiblement inclinés} par rapport à cet axe. Les angles paraxiaux sont petits et s'expriment en radians :
\[\sin\theta\simeq\tan\theta\simeq\theta.\]
Dans ce cadre, une lentille mince réalise un stigmatisme et un aplanétisme approchés. Un petit objet plan perpendiculaire à l'axe a une image approximativement plane et perpendiculaire à l'axe.
\end{loi}
\fig{gauss}{Le diaphragme réduit les rayons éloignés de l'axe ; il ne rend pas tout objet paraxial.}
\Sub{Le rôle du détecteur}
Un point peut former une petite tache et être jugé net si celle-ci est plus petite que le détail résolu par le détecteur. Sur un capteur numérique, on compare sa taille à une dimension de pixel ou à un cercle de confusion choisi. Sur la rétine, la résolution angulaire fournit un autre critère.
\begin{methode}{Vérifier une approximation}
À $\theta=0{,}050\ \mathrm{rad}$, $\tan\theta=0{,}05004$ : l'erreur relative est d'environ $0{,}083\,\%$. À $30^\circ$, l'écart relatif dépasse $10\,\%$. « Petit » doit être rapporté à la précision recherchée.
\end{methode}
\begin{piege}{Conditions et conséquences}
« Les rayons convergent » n'est pas une condition de Gauss. Proximité de l'axe et faible inclinaison sont les conditions ; stigmatisme et aplanétisme approchés en sont les conséquences.
\end{piege}

\Page{Cours\quad /\quad 7}{Lentilles et rayons particuliers}{c-rayons}
Une lentille mince est représentée par un plan passant par son centre optique $O$. L'axe est orienté dans le sens de propagation incidente, ici de gauche à droite. On note $f'=\overline{OF'}$ la distance focale image et $V=1/f'$ la vergence, en $\mathrm{m^{-1}}$ ou dioptries ($\delta$).
\begin{loi}{Trois propriétés de construction}
\begin{itemize}
\item Le rayon passant par $O$ n'est pas dévié, dans le modèle de lentille mince.
\item Un rayon incident parallèle à l'axe ressort en passant par $F'$, ou en semblant en provenir.
\item Un rayon incident passant par $F$, ou dirigé vers $F$, ressort parallèle à l'axe.
\end{itemize}
Dans le même milieu de part et d'autre, $\overline{OF}=-f'$. Une lentille convergente a $f'>0$ ; une lentille divergente a $f'<0$.
\end{loi}
\fig{foyers}{Un faisceau parallèle incliné converge dans le plan focal image en un foyer secondaire.}
\Sub{Foyers secondaires et objet à l'infini}
Chaque direction incidente correspond à un point du plan focal image. Pour le construire, tracer le rayon de cette direction qui traverse $O$ ; son intersection avec le plan passant par $F'$ donne l'image. À une direction faisant un angle algébrique $\theta$ avec l'axe correspond $y'\simeq f'\theta$.
Le \textbf{plan focal objet}, perpendiculaire à l'axe en $F$, est l'ensemble des points dont les faisceaux émergents sont parallèles. Les foyers secondaires sont les points hors axe de ces plans.
\begin{piege}{Le foyer est un point ; le plan focal est un plan}
Un point à l'infini hors de l'axe ne donne pas une image en $F'$. Son image appartient au plan focal image et se trouve hors de l'axe.
\end{piege}

\Page{Cours\quad /\quad 8}{Construire une image réelle}{c-reelle}
\begin{methode}{Une construction entièrement justifiée}
Pour un petit objet $AB$ perpendiculaire à l'axe, placer $A$, $B$, $O$, $F$ et $F'$. Depuis $B$, tracer deux rayons particuliers ; leur intersection donne $B'$. Le point $A'$ est la projection de $B'$ sur l'axe. Ajouter les flèches de propagation et distinguer les prolongements des trajets réels.
\end{methode}
\fig{lentille-reelle}{Objet au-delà de $2f'$ : image réelle, renversée et réduite. Les positions du dessin sont calculées.}
\begin{loi}{Relations de Descartes}
Dans les conditions de Gauss, pour une lentille mince dans le même milieu de part et d'autre :
\[\result{\frac{1}{\oap}-\frac{1}{\oa}=\frac{1}{f'}},\qquad
\result{\gamma=\frac{\overline{A'B'}}{\overline{AB}}=\frac{\oap}{\oa}}.\]
Le grandissement $\gamma$ est algébrique. $\gamma<0$ signifie que l'image est renversée ; $|\gamma|>1$ qu'elle est agrandie.
\end{loi}
\Sub{Exemple avec unités cohérentes}
Une lentille a $f'=15{,}0\ \mathrm{cm}$. L'objet est à $\oa=-40{,}0\ \mathrm{cm}$, de taille $\overline{AB}=2{,}00\ \mathrm{cm}$.
\[\frac1{\oap}=\frac1{15{,}0}-\frac1{40{,}0}=\frac1{24{,}0}\ \mathrm{cm^{-1}}.\]
Ainsi $\oap=+24{,}0\ \mathrm{cm}$, $\gamma=-0{,}600$ et $\overline{A'B'}=-1{,}20\ \mathrm{cm}$. Le calcul et le dessin donnent les mêmes nature et orientation.

\Page{Cours\quad /\quad 9}{Construire une image virtuelle}{c-virtuelle}
\Sub{La loupe : objet entre $F$ et $O$}
\fig{lentille-virtuelle}{Lentille convergente : les rayons émergents divergent ; leurs prolongements se coupent du côté de l'objet.}
Ici $f'=2$ unités et $\oa=-1$ unité : $\oap=-2$ unités et $\gamma=+2$. L'image est virtuelle, droite et agrandie. Les traits pointillés ne représentent pas de la lumière qui repart vers la gauche.
\Sub{Une lentille divergente et un objet réel}
\fig{lentille-divergente}{Lentille divergente : l'image d'un objet réel est virtuelle, droite et réduite.}
Ici $f'=-2$ unités et $\oa=-4$ unités : $\oap=-4/3$ unités et $\gamma=1/3$. Le foyer image $F'$ est à gauche et le foyer objet $F$ à droite. Les propriétés des rayons restent valables avec des prolongements.
\begin{piege}{La nature dépend du faisceau}
« Convergente » qualifie la lentille ; « réelle » qualifie l'image. Une lentille convergente ne donne pas systématiquement une image réelle. Une lentille divergente peut donner une image réelle si elle reçoit un objet virtuel adapté.
\end{piege}

\Page{Cours\quad /\quad 10}{Distances et formules de Newton}{c-newton}
\fig{distances}{Une distance algébrique est positive dans le sens choisi pour l'axe ; l'origine intervient dans sa définition.}
\Sub{Passer de Descartes à Newton}
Posons $p=\oa$, $q=\oap$, $x=\overline{FA}=p+f'$ et $x'=\overline{F'A'}=q-f'$. La relation $1/q-1/p=1/f'$ donne $pq=f'(p-q)$, donc
\[(p+f')(q-f')=pq-f'p+f'q-f'^2=-f'^2.\]
\begin{loi}{Relations de Newton et grandissement}
\[\result{\overline{FA}\,\overline{F'A'}=-f'^2},\qquad
\result{\gamma=\frac{f'}{\overline{FA}}=-\frac{\overline{F'A'}}{f'}}.\]
Les expressions du grandissement découlent de $q/p$ et des mêmes substitutions. Elles s'utilisent avec des distances algébriques, jamais avec des longueurs arbitrairement rendues positives.
\end{loi}
\Sub{Choisir la bonne origine}
Si l'objet est donné à $2{,}0\ \mathrm{mm}$ avant le foyer objet d'une lentille de $f'=20{,}0\ \mathrm{mm}$, on dispose directement de $\overline{FA}=-2{,}0\ \mathrm{mm}$. Newton donne $\overline{F'A'}=+200\ \mathrm{mm}$, puis $\oap=220\ \mathrm{mm}$ et $\gamma=-10$.
\begin{methode}{Contrôle de cohérence en quatre points}
\begin{enumerate}
\item Toutes les distances d'une même équation utilisent la même unité.
\item Les signes correspondent au schéma orienté.
\item Chaque membre de la conjugaison a la dimension d'un inverse de longueur.
\item Un grandissement est sans dimension ; son signe décrit l'orientation.
\end{enumerate}
\end{methode}
Un objet placé exactement en $F$ donne une image à l'infini : la formule impose $1/q=0$. Ce cas ne se traite pas en divisant par une distance nulle dans l'écriture de Newton.

\Page{Cours\quad /\quad 11}{Projeter sur un écran}{c-projection}
On considère un objet réel et un écran, séparés par une distance $D>0$. Une lentille convergente de focale $f'>0$ est placée entre eux. Notons $x>0$ la distance objet-lentille : $\oa=-x$ et $\oap=D-x>0$.
\Sub{Démonstration de la condition de projection}
\[\frac1{D-x}+\frac1x=\frac1{f'}
\quad\Longleftrightarrow\quad x(D-x)=Df'
\quad\Longleftrightarrow\quad x^2-Dx+Df'=0.\]
Le discriminant est $D(D-4f')$. Une position physique existe donc si et seulement si $D\geq4f'$, avec $0<x<D$.
\begin{loi}{Une condition à établir}
\[\result{D\geq4f'}.\]
Si $D>4f'$, deux positions donnent une image nette. Si $D=4f'$, une seule position convient : l'objet et l'image sont à $2f'$ de la lentille et $\gamma=-1$. Si $D<4f'$, aucune position de la lentille entre objet et écran ne convient.
\end{loi}
\fig{bessel}{Les deux positions s'échangent les distances objet et image.}
\Sub{Application : méthode de Bessel}
Pour $D>4f'$, les solutions sont $x_{1,2}=(D\mp\sqrt{D(D-4f')})/2$. Leur écart $d$ vérifie $d^2=D(D-4f')$, d'où
\[\result{f'=\frac{D^2-d^2}{4D}}.\]
Cette formule est une application à retrouver, pas un remplacement des relations de conjugaison. Pour $D=1{,}00\ \mathrm{m}$ et $d=0{,}600\ \mathrm{m}$, on obtient $f'=0{,}160\ \mathrm{m}$ et les positions $x_1=0{,}200\ \mathrm{m}$, $x_2=0{,}800\ \mathrm{m}$.

\Page{Cours\quad /\quad 12}{Le modèle de l'œil}{c-oeil}
Dans le modèle de l'\textbf{œil réduit}, une lentille convergente de vergence variable représente l'effet de la cornée et du cristallin. La rétine est un capteur plan fixe à une distance $\ell$ de l'ordre de $17\ \mathrm{mm}$ de la lentille équivalente. La pupille joue le rôle de diaphragme.
\fig{oeil}{L'accommodation fait varier la vergence ; la rétine reste à la même place.}
\begin{loi}{Accommodation et points limites}
Pour un objet réel à distance $d>0$, $\oa=-d$ et $\oap=\ell$ :
\[\result{V=\frac1\ell+\frac1d}.\]
Le \textbf{punctum remotum} est le point le plus éloigné vu net sans accommodation. Il est à l'infini pour un œil emmétrope. Le \textbf{punctum proximum} est le point le plus proche vu net à accommodation maximale.
\end{loi}
\Sub{Ordres de grandeur à connaître}
Pour un jeune œil emmétrope, on adopte souvent une distance minimale conventionnelle de vision distincte de $25\ \mathrm{cm}$ ; elle varie beaucoup avec l'âge et la personne. La plage de vision nette s'étend alors approximativement de $25\ \mathrm{cm}$ à l'infini.
La limite de résolution angulaire usuelle de l'œil est de l'ordre d'\textbf{une minute d'arc}, soit $\pi/10800\simeq2{,}9\times10^{-4}\ \mathrm{rad}$. Deux détails séparés de $s$ à une distance $d\gg s$ sous-tendent $\theta\simeq s/d$.
\begin{methode}{Exemple}
Avec $\ell=17\ \mathrm{mm}$, la vergence au repos vaut $58{,}8\ \delta$. À $d=25\ \mathrm{cm}$, elle vaut $62{,}8\ \delta$ : l'effort d'accommodation correspond à une variation de $4{,}0\ \delta$.
\end{methode}

\Page{Cours\quad /\quad 13}{L'appareil photographique}{c-photo}
Un appareil photographique est modélisé par une lentille convergente, un diaphragme et un capteur. La mise au point ajuste la distance objectif-capteur ; dans le modèle élémentaire, la focale reste fixe lors de ce réglage.
\fig{photo}{Le capteur reçoit une image réelle et généralement renversée. Le dessin représente le modèle, sans échelle métrique.}
\Sub{Focale et cadrage}
Pour un objet lointain, l'image se forme près du plan focal image. À un petit angle de champ $\theta$ correspond $y'\simeq f'\theta$. Augmenter la focale agrandit l'image d'un objet donné et réduit le champ cadré par un capteur de taille fixée.
Pour un objet réel à distance $p>f'$, la distance de mise au point est
\[\result{q=\frac{f'p}{p-f'}},\qquad |\gamma|=\frac{q}{p}.\]
\begin{loi}{Trois réglages à distinguer}
\textbf{Focale} : agit sur le grandissement et le champ.\par
\textbf{Ouverture} : le nombre d'ouverture est $N=f'/a$, où $a$ est le diamètre du diaphragme. Fermer augmente $N$, diminue la lumière reçue et réduit le flou géométrique de défocalisation.\par
\textbf{Durée d'exposition} : une durée plus grande augmente l'énergie reçue mais aussi le flou de mouvement. À éclairement de scène fixé, le modèle simple donne une exposition proportionnelle à $T/N^2$.
\end{loi}
\begin{piege}{Fermer ne rend pas toujours l'image plus nette}
À très petite ouverture, la diffraction augmente. L'amélioration de la profondeur de champ géométrique doit être distinguée de la résolution globale de l'instrument.
\end{piege}

\Page{Cours\quad /\quad 14}{Construire la profondeur de champ}{c-profondeur}
Le capteur est fixé au plan $q_0$, correspondant à une mise au point sur un objet à distance $p_0$. Un autre point de l'axe, à distance $p$, aurait son image ponctuelle au plan $q=f'p/(p-f')$. Sur le capteur, le faisceau dessine donc un disque plutôt qu'un point.
\fig{profondeur}{Suivre les deux rayons issus des bords du diaphragme. Leur écart au plan du capteur donne le diamètre de flou.}
\Sub{Construction géométrique exigible}
\begin{enumerate}
\item Construire l'image idéale du point objet en utilisant la conjugaison.
\item Relier les bords de l'ouverture à ce point image.
\item Intersecter le cône ainsi obtenu avec le plan du capteur.
\item Déclarer le point acceptable si le diamètre $b$ de la tache est inférieur au diamètre admissible $\varepsilon$.
\end{enumerate}
\begin{loi}{Expression obtenue par triangles semblables}
Si le diaphragme de diamètre $a$ est dans le plan de la lentille :
\[\result{b=a\left|1-\frac{q_0}{q}\right|},\qquad
\text{image acceptable si }b\leq\varepsilon.\]
La \textbf{profondeur de champ} est l'intervalle des distances objets satisfaisant ce critère pour le réglage choisi. Elle dépend donc du critère de netteté, et pas seulement de la lentille.
\end{loi}
Fermer le diaphragme diminue $a$ et élargit l'intervalle acceptable dans ce modèle. Pour chiffrer un effet de focale, il faut préciser ce qui reste fixé : distance de mise au point, cadrage, taille du capteur et nombre d'ouverture.
\source{Le défi D1 déduit les bornes de cet intervalle sans demander de mémoriser une formule d'hyperfocale.}

\Page{Cours\quad /\quad 15}{Guider dans une fibre optique}{c-fibre}
Une fibre à saut d'indice possède un cœur d'indice $n_1$ entouré d'une gaine d'indice $n_2<n_1$. On étudie les rayons \textbf{méridiens}, contenus dans un plan passant par l'axe. Le milieu d'entrée a l'indice $n_0$ et la face d'entrée est perpendiculaire à l'axe.
\fig{fibre}{Attention : l'incidence sur la paroi se mesure à la normale locale, et non à l'axe de la fibre.}
\Sub{Démonstration du cône d'acceptance}
À l'entrée, $n_0\sin\alpha=n_1\sin\theta$. À la paroi cœur-gaine, l'incidence vaut $i=\pi/2-\theta$. Pour la réflexion totale, $\sin i>n_2/n_1$, donc $\cos\theta>n_2/n_1$. Ainsi
\[n_0^2\sin^2\alpha=n_1^2\sin^2\theta<n_1^2-n_2^2.\]
\begin{loi}{Ouverture numérique et angle limite}
\[\mathrm{ON}=\sqrt{n_1^2-n_2^2},\qquad
\result{\alpha_{\max}=\arcsin\!\left(\frac{\mathrm{ON}}{n_0}\right)}\]
si $\mathrm{ON}\leq n_0$. Les rayons effectivement guidés par réflexion totale satisfont $|\alpha|<\alpha_{\max}$ ; l'égalité décrit le seuil. Si $\mathrm{ON}>n_0$, tous les angles d'entrée physiquement accessibles satisfont la condition géométrique.
\end{loi}
\Sub{Application et limites}
Avec $n_1=1{,}48$, $n_2=1{,}46$ et $n_0=1{,}00$, $\mathrm{ON}=0{,}242$ et $\alpha_{\max}=14{,}0^\circ$. Le cône d'entrée complet a une ouverture angulaire de $2\alpha_{\max}$.
Le modèle néglige les pertes, les courbures et la structure ondulatoire des modes. Il est adapté à l'interprétation géométrique d'une fibre multimode de dimensions grandes devant $\lambda$.

\Page{Cours\quad /\quad 16}{La dispersion intermodale}{c-dispersion}
Deux rayons guidés peuvent parcourir des longueurs différentes pour la même longueur axiale $L$. Une impulsion brève se répartit donc sur plusieurs temps d'arrivée : c'est la \textbf{dispersion intermodale}. On suppose les indices indépendants de la fréquence dans le calcul qui suit.
\fig{dispersion}{Les trajets axial et oblique n'ont pas la même durée.}
\Sub{Démonstration à savoir reconstruire}
Pour un rayon d'angle interne $\theta$, la longueur parcourue est $L/\cos\theta$ ; avec la célérité $c/n_1$, on obtient
\[t(\theta)=\frac{n_1L}{c\cos\theta}.
\]
Le rayon axial est le plus rapide : $t_{\min}=n_1L/c$. Au seuil de guidage, $\cos\theta_{\max}=n_2/n_1$, donc la borne supérieure des temps vaut $t_{\max}=n_1^2L/(cn_2)$.
\begin{loi}{Étalement temporel maximal géométrique}
\[\result{\Delta t=\frac{n_1L}{c}\left(\frac{n_1}{n_2}-1\right)}.\]
La valeur limite est un supremum pour les rayons strictement guidés. Si $n_1-n_2\ll n_1$, alors $\Delta t\simeq L(n_1-n_2)/c$.
\end{loi}
\Sub{Application}
Avec $L=1{,}00\ \mathrm{km}$, $n_1=1{,}48$, $n_2=1{,}46$ et $c=3{,}00\times10^8\ \mathrm{m\,s^{-1}}$ : $t_{\min}=4{,}93\ \mu\mathrm{s}$ et $\Delta t=67{,}6\ \mathrm{ns}$. Une largeur d'impulsion initiale non nulle s'ajoute au problème.
\begin{piege}{Un critère de débit doit être annoncé}
Si l'on impose une séparation des impulsions de $2\Delta t$, un débit binaire maximal estimé est $R\leq1/(2\Delta t)$. Le facteur 2 relève de ce critère choisi ; ce n'est pas une loi universelle de la fibre.
\end{piege}

\Page{Pratique\quad /\quad TP}{Former et mesurer une image}{c-tp}
\textbf{Objectif.} Mesurer la focale d'une lentille convergente, confronter plusieurs méthodes et présenter un résultat avec une incertitude. Matériel : banc gradué, objet diffus éclairé par LED, lentille convergente, écran et règle. Ne pas utiliser le Soleil comme objet ni regarder dans un faisceau laser.
\fig{banc}{Mesurer les positions des plans utiles : objet, centre optique approché et écran.}
\Sub{Protocole à concevoir puis à réaliser}
\begin{enumerate}
\item Aligner objet, lentille et écran à la même hauteur. Réaliser une image nette et vérifier qu'elle est renversée.
\item Mesurer $x_A$, $x_L$, $x_E$ puis calculer $p=x_L-x_A>0$ et $q=x_E-x_L>0$. Estimer l'intervalle de positions de l'écran encore jugées nettes.
\item Refaire au moins cinq mises au point indépendantes, avec plusieurs distances objets. Calculer chaque focale par $f'=pq/(p+q)$.
\item Fixer une distance $D>4f'$ ; rechercher les deux positions de Bessel et mesurer leur écart $d$. Comparer les résultats.
\item Expliquer l'effet d'un diaphragme sur la luminosité, la netteté perçue et la précision de mise au point.
\end{enumerate}
\begin{docu}{Jeu d'essai simulé, pour préparer l'exploitation}
Ces valeurs ne sont pas des mesures réellement réalisées.
\begin{center}\begin{tabular}{c r r r}
\toprule Essai & 1 & 2 & 3\\\midrule
$p$ (cm)&30,0&40,0&60,0\\
$q$ (cm)&30,1&24,0&20,0\\\bottomrule
\end{tabular}\end{center}
Les focales calculées valent environ $15{,}025$, $15{,}000$ et $15{,}000\ \mathrm{cm}$. L'accord des valeurs ne suffit pas à établir à lui seul l'incertitude instrumentale.
\end{docu}
\source{Compte rendu attendu : schéma légendé, procédure, tableau avec unités, incertitudes, résultat, comparaison et limites du modèle.}

\Page{Pratique\quad /\quad numérique}{Incertitude et simulation}{c-numerique}
Pour $f'=pq/(p+q)$ et des valeurs $p,q$ supposées indépendantes, la propagation au premier ordre donne
\[u(f')=\sqrt{\left(\frac{q^2}{(p+q)^2}u(p)\right)^2+
\left(\frac{p^2}{(p+q)^2}u(q)\right)^2}.\]
Pour $p=40{,}0\ \mathrm{cm}$, $q=24{,}0\ \mathrm{cm}$ et $u(p)=u(q)=0{,}10\ \mathrm{cm}$, on obtient $u(f')=0{,}0415\ \mathrm{cm}$. On peut écrire $f'=15{,}00\ \mathrm{cm}$ avec une incertitude-type $u(f')=0{,}04\ \mathrm{cm}$.
\begin{piege}{Attention à une position commune}
Sur un banc, $p=x_L-x_A$ et $q=x_E-x_L$ partagent $x_L$. Si les trois positions sont indépendantes, $\operatorname{cov}(p,q)=-u(x_L)^2$. Les traiter indépendamment serait alors injustifié. Simuler directement les trois positions évite cet oubli.
\end{piege}
\Sub{Une simulation reproductible en Python}
\begingroup\small
\begin{verbatim}
import numpy as np
rng = np.random.default_rng(2026)
M = 100_000
# Toutes les positions sont en centimetres.
# Hypothese : erreurs normales independantes,
# d'incertitude-type 0.10 cm sur chaque position.
xA = rng.normal(0.0, 0.10, M)
xL = rng.normal(40.0, 0.10, M)
xE = rng.normal(64.0, 0.10, M)
p, q = xL-xA, xE-xL
f = p*q/(p+q)
print(np.mean(f), np.std(f, ddof=1))
\end{verbatim}
\endgroup
Le résultat attendu est proche de $15{,}00\ \mathrm{cm}$ avec une dispersion de $0{,}0485\ \mathrm{cm}$ ; les derniers chiffres fluctuent avec le tirage. Les lois normales sont des hypothèses de simulation, pas une propriété automatique des règles graduées.
\begin{methode}{Comparer deux résultats}
Pour des résultats indépendants $f_1,f_2$, calculer $z=|f_1-f_2|/\sqrt{u(f_1)^2+u(f_2)^2}$. Une valeur inférieure à 2 est un repère usuel de compatibilité ; elle ne prouve ni l'exactitude du modèle ni l'absence de biais commun.
\end{methode}

\Page{Bilan de cours}{Les réflexes à conserver}{c-synthese}
\begin{loi}{Les relations et leurs conditions}
\begin{align*}
&\lambda_0=c/\nu,\quad n=c/v,\quad\lambda=\lambda_0/n,\\
&n_1\sin i=n_2\sin r,\quad i_c=\arcsin(n_2/n_1)\quad(n_1>n_2),\\
&\frac1{\oap}-\frac1{\oa}=\frac1{f'},\quad
\gamma=\frac{\oap}{\oa},\quad V=1/f',\\
&\overline{FA}\,\overline{F'A'}=-f'^2,\quad
\gamma=\frac{f'}{\overline{FA}}=-\frac{\overline{F'A'}}{f'},\\
&D\geq4f'\quad\text{pour projeter avec une lentille entre objet et écran},\\
&n_0\sin\alpha_{\max}=\sqrt{n_1^2-n_2^2},\\
&\Delta t=\frac{n_1L}{c}\left(\frac{n_1}{n_2}-1\right).
\end{align*}
Les relations des lentilles demandent le modèle mince, les conditions de Gauss et le même milieu de chaque côté. Celles de la fibre demandent le modèle géométrique méridien sans dispersion matérielle.
\end{loi}
\Sub{Cinq démonstrations à refaire sans corrigé}
\begin{enumerate}
\item La condition de réflexion totale à partir de l'existence de $r$.
\item Les relations de Newton à partir de Descartes.
\item La condition $D\geq4f'$ et les deux positions de projection.
\item Le cône d'acceptance à partir des angles d'entrée et de paroi.
\item La dispersion intermodale à partir des longueurs parcourues.
\end{enumerate}
\Sub{Avant de rendre une copie}
\checkline{Mon schéma comporte axe orienté, normales, foyers, noms et flèches.}
\checkline{Mes pointillés désignent des prolongements, pas des trajets réels.}
\checkline{J'annonce les hypothèses et je distingue longueurs et distances algébriques.}
\checkline{Je donne une expression littérale, une unité et une interprétation.}
\checkline{Je vérifie l'ordre de grandeur, la nature et l'orientation de l'image.}
